\chapter{Kommentierte Quellen}

Achtundzwanzig Werke, eingeordnet: Zu jedem steht, wofür man es aufschlägt und was es leistet, das andere nicht leisten.

\lit{1}{Angrist, J., Pischke, J.-S.: \textit{Mostly Harmless Econometrics}. Princeton 2009.}{\textbf{Das Buch der Glaubwürdigkeitswende und die Hauptquelle dieses Bandes für die Kapitel~3, 7 und 9.} Instrumente, feste Effekte, doppelte Differenz, Schwellen --- alles aus der Sicht des Entwurfs, mit Beispielen aus Arbeitsmarkt und Bildung. Knapp, meinungsstark, unverzichtbar.}

\lit{2}{Angrist, J., Pischke, J.-S.: \textit{The Credibility Revolution in Empirical Economics}. J.\ Econ.\ Perspect.\ 24, 2010.}{Der Aufsatz, der der Wende ihren Namen gab --- und die Gegenrede von Leamer, Sims und Keane im selben Heft. \textbf{Wer den Streit aus Kapitel~15 verstehen will, lese alle vier Beiträge.}}

\lit{3}{Arellano, M., Bond, S.: \textit{Some Tests of Specification for Panel Data}. Rev.\ Econ.\ Stud.\ 58, 1991.}{Der GMM-Schätzer für dynamische Panels aus Kapitel~8, mit dem Test auf serielle Korrelation. Einer der meistzitierten Aufsätze der Ökonometrie; zusammen mit Roodman [25] zu lesen.}

\lit{4}{Auer, L.\ von: \textit{Ökonometrie --- Eine Einführung}. 8.\ Aufl., Springer Gabler 2023.}{\textbf{Auf Deutsch, das verbreitetste Lehrbuch im deutschen Sprachraum.} Regression, Verletzungen der Annahmen, Instrumente, Zeitreihen --- schrittweise und mit vielen Aufgaben. Der Einstieg für Leser, die Band 17 kennen und die Kapitel~1 bis 3 und 10 bis 11 vertiefen wollen.}

\lit{5}{Cameron, A.\,C., Trivedi, P.: \textit{Microeconometrics: Methods and Applications}. Cambridge 2005.}{Das Nachschlagewerk zu den Teilen~II und III: GMM, diskrete Wahl, Zensierung, Auswahl, Panels --- vollständig und mit Beispielen aus der Gesundheitsökonomie, darunter das Zweiteilmodell. Über tausend Seiten; die Fundstelle für alles, was hier nur angerissen wurde.}

\lit{6}{Card, D., Krueger, A.: \textit{Minimum Wages and Employment}. Amer.\ Econ.\ Rev.\ 84, 1994.}{Die berühmteste doppelte Differenz --- Schnellrestaurants in New Jersey und Pennsylvania. \textbf{Lesenswert als Muster für den Entwurf}: wie man eine Kontrollgruppe begründet, prüft und verteidigt.}

\lit{7}{Cunningham, S.: \textit{Causal Inference: The Mixtape}. Yale 2021.}{Die moderne Darstellung der Kapitel~3, 7 und 9 einschließlich der gestaffelten Differenz aus dem Bruchfall von Kapitel~9. Frei verfügbar, mit Programmcode, unterhaltsam geschrieben --- die Ergänzung zu [1] um das, was seit 2009 dazukam.}

\lit{8}{Engle, R., Granger, C.: \textit{Co-Integration and Error Correction}. Econometrica 55, 1987.}{Der Aufsatz zu Kapitel~12 --- Definition, Test, Fehlerkorrektur, Darstellungssatz. Schwer zu lesen und einer der folgenreichsten der Zeitreihenökonometrie; der Nobelpreis von 2003 beruhte vor allem auf ihm.}

\lit{9}{Granger, C., Newbold, P.: \textit{Spurious Regressions in Econometrics}. J.\ Econometrics 2, 1974.}{Das Experiment aus Kapitel~11 im Original, auf acht Seiten. \textbf{Wer es liest, versteht, warum ein Fach seine Praxis ändern musste} --- und warum es fünfzig Jahre brauchte, bis es Yules Warnung von 1926 hörte.}

\lit{10}{Greene, W.: \textit{Econometric Analysis}. 8.\ Aufl., Pearson 2018.}{Das umfassendste Lehrbuch des Fachs --- Rangbedingung, Simultanität, GMM, diskrete Wahl, Panels, Zeitreihen. Die Fundstelle für die erste Auslassung. Als Nachschlagewerk unübertroffen, zum Lesen zu dicht.}

\lit{11}{Hamilton, J.: \textit{Time Series Analysis}. Princeton 1994.}{\textbf{Das Standardwerk zu Teil~IV}: ARMA, Einheitswurzeln, Kointegration, VAR, Kalman-Filter, mit allen Beweisen einschließlich der Brownschen Grenzverteilungen aus Kapitel~11 und des Darstellungssatzes aus Kapitel~12. Dreißig Jahre alt und nicht ersetzt.}

\lit{12}{Hayashi, F.: \textit{Econometrics}. Princeton 2000.}{Die Ökonometrie konsequent als GMM aufgebaut --- jeder Schätzer ein Sonderfall von Kapitel~4. Streng, elegant, mit vollständigen Beweisen der Konsistenz und Effizienz. Die Fundstelle für die zweite Auslassung und das Buch für Leser, die die Struktur hinter den Verfahren sehen wollen.}

\lit{13}{Heckman, J.: \textit{Sample Selection Bias as a Specification Error}. Econometrica 47, 1979.}{Satz 6.2 im Original. Kurz und klar; der Titel ist die These. Zusammen mit Heckmans späteren Warnungen vor der mechanischen Anwendung zu lesen, etwa in seiner Nobelpreisrede von 2000.}

\lit{14}{Hübler, O.: \textit{Ökonometrie}. UTB 2005.}{\textbf{Auf Deutsch, mit Schwerpunkt auf Mikroökonometrie}: diskrete Wahl, Zensierung, Auswahl, Panels --- die Kapitel~5 bis 8 dieses Bandes. Knapper als [5] und für deutschsprachige Leser der direktere Weg.}

\lit{15}{Imbens, G., Rubin, D.: \textit{Causal Inference for Statistics, Social, and Biomedical Sciences}. Cambridge 2015.}{Die potentiellen Ergebnisse aus Band 23 in voller Breite, einschließlich der Deutung von Instrumenten als Befolgereffekt. \textbf{Die Brücke zwischen diesem Band und Band 23} --- Statistiker und Ökonometriker in einem Buch.}

\lit{16}{Kennedy, P.: \textit{A Guide to Econometrics}. 6.\ Aufl., Wiley-Blackwell 2008.}{\textbf{Ökonometrie ohne Formeln --- oder fast.} Jedes Kapitel erklärt zuerst die Idee, dann die Technik, dann die Fallstricke. Für Leser, die verstehen wollen, warum man etwas tut, bevor sie lernen, wie. Enthält ein berühmtes Kapitel über die zehn Gebote der angewandten Ökonometrie.}

\lit{17}{Lucas, R.: \textit{Econometric Policy Evaluation: A Critique}. Carnegie-Rochester Conf.\ 1, 1976.}{Die Kritik aus Kapitel~13 im Original, mit den drei Beispielen Konsum, Investition und Phillips-Kurve. \textbf{Einer der einflussreichsten Aufsätze der Wirtschaftswissenschaften}, und kürzer, als sein Ruf vermuten lässt.}

\lit{18}{Makridakis, S., Spiliotis, E., Assimakopoulos, V.: \textit{The M4 Competition}. Int.\ J.\ Forecasting 36, 2020.}{Die Ergebnisse des vierten Wettbewerbs und die Lehren aus vierzig Jahren --- Abschnitt 13.3. Nüchtern und für jeden lesbar, der Vorhersagen beurteilen muss.}

\lit{19}{Manski, C.: \textit{Identification for Prediction and Decision}. Harvard 2007.}{Die Lehre von den Schranken aus den Forschungskästen der Kapitel~2 und 6: Was folgt, wenn man die unhaltbare Annahme weglässt? \textbf{Die methodische Gegenposition zur Punktschätzung}, streng und mit ungewöhnlicher Klarheit über das, was man nicht weiß.}

\lit{20}{Newhouse, J.\ und die Insurance Experiment Group: \textit{Free for All? Lessons from the RAND Health Insurance Experiment}. Harvard 1993.}{\textbf{Der Leitfall dieses Bandes, von denen beschrieben, die ihn durchgeführt haben.} Entwurf, Durchführung, Ergebnisse, Einschränkungen --- und das Zweiteilmodell, das dafür entwickelt wurde. Für die Gesundheitsökonomie, was Fishers Feldversuche für die Agrarwissenschaft sind.}

\lit{21}{Roth, J., Sant'Anna, P., Bilinski, A., Poe, J.: \textit{What's Trending in Difference-in-Differences?} J.\ Econometrics 235, 2023.}{Die Übersicht über die Krise und Rettung der doppelten Differenz --- Bruchfall und Forschungskasten von Kapitel~9. \textbf{Wer eine gestaffelte Einführung auswertet, muss das gelesen haben}; es ist die zugänglichste Zusammenfassung einer unübersichtlichen Literatur.}

\lit{22}{Schira, J.: \textit{Statistische Methoden der VWL und BWL}. 6.\ Aufl., Pearson 2021.}{\textbf{Auf Deutsch, als Brücke von Band 17 zu diesem Band.} Regression, Zeitreihen, Grundbegriffe der Ökonometrie, mit wirtschaftswissenschaftlichen Beispielen. Für Leser, die zuerst die Statistik wiederholen wollen, bevor sie die Endogenität angehen.}

\lit{23}{Sims, C.: \textit{Macroeconomics and Reality}. Econometrica 48, 1980.}{Die Kritik an den Ausschlüssen aus Kapitel~2 und der Vorschlag der Vektorautoregression aus Kapitel~12. \textbf{Der Gegenentwurf zu Lucas [17]} --- beide Aufsätze zusammen umreißen den Streit, der die Makroökonometrie seither prägt.}

\lit{24}{Stock, J., Watson, M.: \textit{Introduction to Econometrics}. 4.\ Aufl., Pearson 2019.}{Das verbreitetste englischsprachige Lehrbuch --- \textbf{der beste Einstieg in das ganze Fach}, mit Instrumenten, Panels, Zeitreihen und einem Kapitel über schwache Instrumente, das Stock selbst mitbegründet hat. Weniger streng als [10], klarer als alle anderen.}

\lit{25}{Roodman, D.: \textit{A Note on the Theme of Too Many Instruments}. Oxf.\ Bull.\ Econ.\ Stat.\ 71, 2009.}{Die Warnung aus Kapitel~8: Was geschieht, wenn man das Verfahren von [3] mit Dutzenden Instrumenten füttert. Kurz, praktisch, mit Faustregeln --- und der Nachweis, dass der Hansen-Test dann nichts mehr prüft.}

\lit{26}{Winker, P.: \textit{Empirische Wirtschaftsforschung und Ökonometrie}. 4.\ Aufl., Springer 2017.}{\textbf{Auf Deutsch, mit Blick auf die Praxis}: Datenquellen, Spezifikation, Diagnostik, Berichtsweise --- die Kapitel~14 und 15 dieses Bandes in Lehrbuchform. Für Leser, die Studien lesen und beurteilen müssen, mehr als für solche, die sie rechnen.}

\lit{27}{Wooldridge, J.: \textit{Econometric Analysis of Cross Section and Panel Data}. 2.\ Aufl., MIT Press 2010.}{\textbf{Das maßgebliche Werk zu den Teilen~I bis III} --- Endogenität, GMM, diskrete Wahl, Zensierung, Auswahl, Panels, alles mit Beweisen und in einheitlicher Notation. Die Fundstelle für die Feinheiten der Kapitel~5 bis 8, die hier übergangen wurden.}

\lit{28}{Wooldridge, J.: \textit{Introductory Econometrics: A Modern Approach}. 7.\ Aufl., Cengage 2020.}{Die zugängliche Fassung von [27]. Beginnt bei der Regression und endet bei Panels und Zeitreihen; jedes Kapitel mit Datensätzen zum Nachrechnen. \textbf{Das Buch, mit dem die meisten Ökonometriker begonnen haben} --- und für Leser dieses Bandes die beste Übungssammlung.}

\begin{merksatz}
\boxtitel{Wenn Sie nur drei Werke lesen}
\begin{center}
\renewcommand{\arraystretch}{1.6}
\begin{tabular}{@{}P{1.0cm}P{3.4cm}P{3.8cm}@{}}
\toprule
\textbf{[24]} & \textbf{Stock, Watson} & \textbf{Für alles}\\
\textbf{[1]} & \textbf{Angrist, Pischke} & Für den Entwurf\\
\textbf{[11]} & \textbf{Hamilton} & \textbf{Für die Zeitreihen}\\
\bottomrule
\end{tabular}
\end{center}

\begin{center}
\begin{minipage}{0.88\textwidth}
\kursiv{Stock und Watson decken den ganzen Band ab und sind lesbar. Angrist und Pischke lehren, worauf es bei Querschnitts- und Paneldaten ankommt: den Vergleich, nicht das Verfahren. \textbf{Und Hamilton ist das Buch, das man aufschlägt, wenn eine Zeitreihe mehr verlangt als Kapitel~10 bis 12 geben können.}}
\end{minipage}
\end{center}

\begin{center}
\begin{minipage}{0.88\textwidth}
\kursiv{Nach Interessenlage: [27] Wooldridge für die Theorie der Mikroökonometrie, [5] Cameron und Trivedi für die Gesundheitsökonomie, [7] Cunningham und [21] Roth u.\,a.\ für den Stand der doppelten Differenz, [19] Manski für die Schranken, [20] Newhouse für den Leitfall. Auf Deutsch: [4] von Auer als Lehrbuch, [14] Hübler für die Mikroökonometrie, [22] Schira als Brücke von Band 17, [26] Winker für die Praxis.}
\end{minipage}
\end{center}
\end{merksatz}

\bigskip
\begin{flushright}
\itshape\color{grau}
Die Daten sagen, was geschah.\\
Die Annahme sagt, warum.\\
Erst beides zusammen sagt,\\
was gälte, wenn man handelte.
\end{flushright}
